Ce stage a été une expérience enrichissante et formatrice pour moi. En entrant dans ce projet, je n'avais que peu de connaissances en modélisation océanographique et en programmation. Cependant, grâce à l'encadrement, aux recherches et aux défis rencontrés, j'ai acquis une solide compétence dans ces domaines. Mes codes, autrefois rudimentaires, ont évolué pour devenir plus structurés et professionnels.

J'ai eu la chance de rencontrer des experts internationaux en programmation et en modélisation océanographique, ce qui m'a permis d'élargir mon réseau professionnel et d'acquérir une perspective globale sur le domaine. La charge de travail, bien que intimidante au départ, m'a enseigné une leçon précieuse : avec détermination et volonté, rien n'est insurmontable.

Je suis particulièrement fier d'avoir contribué à QcrocoFlow. J'espère que cet outil facilitera l'accès à la modélisation océanographique CROCO pour de nombreux chercheurs et aidera à résoudre de nouvelles problématiques. Au total, j'ai développé 3590 lignes de code réparties sur 31 scripts, couvrant à la fois l'interface utilisateur et le processus de génération des fichiers d'entrée du modèle.

\subsection{Perspectives}

L'avenir de QcrocoFlow est prometteur. Pour les futurs développeurs ou stagiaires travaillant sur ce projet, plusieurs améliorations peuvent être envisagées :

\begin{itemize}
	\item \textbf{Amélioration de l'interface utilisateur :} Rendre l'interface plus intuitive et ergonomique pour les utilisateurs.
	\item \textbf{Nested Grid :} Cette fonctionnalité permettrait de générer des grilles filles et des grilles parents pour le modèle, offrant ainsi une résolution variable sur différentes zones d'intérêt\footnote{Les grilles "nested" ou imbriquées permettent d'avoir une résolution plus élevée dans des zones d'intérêt spécifiques tout en conservant une résolution plus faible ailleurs, optimisant ainsi les ressources de calcul.}.
	\item \textbf{Édition du masque :} Ajouter une fonction pour permettre aux utilisateurs d'éditer le masque lors de la génération de la grille, offrant ainsi plus de flexibilité.
	\item \textbf{Configuration du modèle CROCO :} La configuration du modèle CROCO est un processus complexe. Intégrer une section dédiée dans QcrocoFlow pour faciliter cette étape serait d'une grande valeur ajoutée.
\end{itemize}

En conclusion, ce stage a été une étape cruciale dans mon parcours professionnel. Je suis reconnaissant pour toutes les opportunités d'apprentissage et les défis qui m'ont été présentés. J'espère que mes contributions à QcrocoFlow bénéficieront à la communauté de modélisation océanographique pendant de nombreuses années.

